\documentclass[11pt]{article}

% ---------- Packages ----------
\usepackage[margin=1in]{geometry}
\usepackage{setspace}
\usepackage{booktabs}
\usepackage{dcolumn}
\usepackage{graphicx}
\usepackage{float}
\usepackage{caption}
\usepackage{subcaption}
\usepackage{amsmath}
\usepackage{hyperref}
\usepackage{natbib}
\usepackage{url}
\usepackage{microtype}
\usepackage{lmodern}
\usepackage{siunitx}
\usepackage{threeparttable}
\usepackage{pdflscape}
\usepackage{xcolor}
\usepackage{appendix}

\onehalfspacing

\newcommand{\agency}{Orange County Sheriff's Department}
\newcommand{\yr}{2020}
\newcommand{\N}[1]{\num{#1}}

\title{\textbf{Racial Disparities in Traffic Stops:\\
\agency{} (\yr{})}\\[0.4em]
\large RIPA Stop Data Analysis}
\author{Leo Wiemelt\\
\small ECON 185}
\date{\today}

\begin{document}

\maketitle
\thispagestyle{empty}

\begin{abstract}
\noindent
This paper studies racial disparities in traffic stops conducted by the
Orange County Sheriff's Department in 2020 using data collected under the
California Racial and Identity Profiling Act (RIPA). I first document racial
differences in stop volumes, stop reasons, and search rates. I then apply the
Veil of Darkness test to examine whether minority drivers are more likely to be
stopped when officers can observe race prior to the stop. Finally, I conduct
an outcome (hit-rate) test to assess whether searches of minority drivers are
less productive. The combined evidence suggests meaningful racial disparities
in policing outcomes, though each test has limitations that preclude a fully
causal interpretation.
\end{abstract}

\bigskip

% ============================================================
\section{Context and Summary Statistics}
\label{sec:summary}
% ============================================================

\subsection*{Data}

I use \yr{} RIPA stop data from the \agency{}.\footnote{\url{https://openjustice.doj.ca.gov/data}}
Orange County is a large, suburban county with a mix of incorporated cities and unincorporated areas policed by the Sheriff’s Department.

The raw dataset includes all reported traffic stops conducted during 2020. I restrict the sample to traffic stops with real values for perceived race, and valid stop timestamps. These restrictions ensure comparability across racial groups and consistency across analyses.

\subsection*{Task A}

Table~\ref{tab:summary} reports total stops by race, the racial composition of
all stops, and a few interesting reasons for stopping: moving violations,
equipment violations, and investigative stops. These categories differ in the
degree of officer discretion involved.

\begin{table}[H]
  \centering
  \caption{Stop Characteristics by Race, \agency{} (\yr{})}
  \label{tab:summary}
  \begin{threeparttable}
    \begin{tabular}{lS[table-format=6.0]S[table-format=2.1]
                    S[table-format=2.1]S[table-format=2.1]S[table-format=2.1]}
      \toprule
      & {Total Stops} & {\% of Stops} & {Moving (\%)} & {Equipment (\%)} & {Investigative (\%)} \\
      \midrule
      White          &     &     &     &     &     \\
      Black          &     &     &     &     &     \\
      Hispanic       &     &     &     &     &     \\
      Asian          &     &     &     &     &     \\
      Other/Unknown  &     &     &     &     &     \\
      \midrule
      \textbf{Total} &     & 100.0 &     &     &     \\
      \bottomrule
    \end{tabular}
    \begin{tablenotes}[flushleft]
      \footnotesize
      \item \textit{Notes:} Percentages in columns 3–5 are computed within race.
      Stop reasons correspond to RIPA variable \texttt{STOP\_REASON\_1}.
      Sample excludes stops with missing race or invalid timestamps.
    \end{tablenotes}
  \end{threeparttable}
\end{table}

\subsection*{Task B}

I define a search indicator equal to one if any search occurred during the stop,
including consent, probable cause, incident-to-arrest, or warrant searches.
This ``any search'' measure captures overall search behavior without
conditioning on search type. Figure~\ref{fig:search_rates} plots search rates by
race with 95\% confidence intervals.

\begin{figure}[H]
  \centering
  \includegraphics[width=0.72\textwidth]{figures/search_rate_by_race.pdf}
  \caption{Search Rate by Race, \agency{} (\yr{})}
  \label{fig:search_rates}
  \footnotesize
  \textit{Notes:} Search rate equals the number of stops involving any search
  divided by total stops. Bars show 95\% confidence intervals.
\end{figure}

\subsection*{Conceptual Question: Does a Disparity Prove Discrimination?}

Differences in search rates across racial groups are consistent with racial
discrimination but do not, on their own, establish discriminatory intent.
Officers may observe factors correlated with race—such as location or driving
behavior—that are not recorded in the data.

A further complication is the benchmark problem: the racial distribution of
drivers on the road is unobserved. Without knowing the appropriate comparison
population, raw disparities cannot be interpreted causally.

These limitations motivate quasi-experimental approaches that attempt to hold
officer information constant. The next sections apply the Veil of Darkness test
and an outcome test to address these concerns.

% ============================================================
\section{The Veil of Darkness Test}
\label{sec:vod}
% ============================================================

\subsection*{Design}

The Veil of Darkness test exploits seasonal variation in sunset times. During
daylight, officers can typically perceive a driver’s race prior to initiating
a stop, whereas after dark this information is largely unavailable. If minority
drivers are stopped at higher rates in daylight than in darkness, conditional
on time of day, this pattern is consistent with racial profiling.

I restrict attention to an intertwilight window around sunset and include
15-minute time-of-day fixed effects. Sunset times are computed using the
\texttt{astral} package at the geographic centroid of Orange County.

\subsection*{Sample Restriction}

The VOD sample includes only stops occurring between 5:00 PM and 8:30 PM. This
window ensures overlap between daylight and darkness across seasons while
limiting changes in traffic composition.

\subsection*{Specification}

Equation~\eqref{eq:vod} estimates the effect of daylight on the probability that
a stopped driver is perceived as Black or Hispanic.

\begin{table}[H]
  \centering
  \caption{Veil-of-Darkness Regression: Effect of Daylight on Minority Stop Rate}
  \label{tab:vod}
  \begin{threeparttable}
    \begin{tabular}{lD{.}{.}{2,4}D{.}{.}{2,4}D{.}{.}{2,4}}
      \toprule
      & (1) & (2) & (3) \\
      & Black or Hispanic & Black & Hispanic \\
      \midrule
      Daylight      &   &   &   \\
                    & ( ) & ( ) & ( ) \\[4pt]
      \midrule
      15-min dummies & Yes & Yes & Yes \\
      Observations  &   &   &   \\
      Mean dep.\ var.\ (dark) & & & \\
      \bottomrule
    \end{tabular}
    \begin{tablenotes}[flushleft]
      \footnotesize
      \item \textit{Notes:} Linear probability model with heteroskedasticity-robust
      standard errors. Sample restricted to the intertwilight window.
    \end{tablenotes}
  \end{threeparttable}
\end{table}

\subsection*{Conceptual Question: Interpreting the VOD Coefficient}

A positive daylight coefficient indicates that minority drivers constitute a
larger share of stops when officers can observe race. This pattern is
consistent with racial profiling but does not identify individual officer
behavior.

The validity of the test relies on the assumption that the racial composition
of drivers does not change systematically between daylight and darkness within
the window. Violations of this assumption could bias the estimate.

% ============================================================
\section{The Outcome Test: Hit Rates}
\label{sec:hitrate}
% ============================================================

\subsection*{Design}

The outcome test compares the success rate of searches across racial groups.
If officers apply lower suspicion thresholds to minority drivers, searches of
those drivers should be less likely to uncover contraband. I define ``any
contraband'' as the discovery of drugs, weapons, or other illegal items.

\subsection*{Specification}

Equation~\eqref{eq:hit} is estimated on the subsample of stops involving a
search, with White drivers as the omitted category.

\begin{table}[H]
  \centering
  \caption{Hit-Rate Regression and Search/Contraband Rates by Race}
  \label{tab:hitrate}
  \begin{threeparttable}
    \begin{tabular}{lD{.}{.}{1,4}S[table-format=2.1]S[table-format=2.1]}
      \toprule
      & (1) & \multicolumn{2}{c}{Raw rates (\%)} \\
      \cmidrule(lr){3-4}
      & Contraband found & Search rate & Hit rate \\
      \midrule
      White (base) & & & \\
      Black & & & \\
            & ( ) & & \\
      Hispanic & & & \\
               & ( ) & & \\
      Asian & & & \\
            & ( ) & & \\
      \midrule
      Observations & & & \\
      \bottomrule
    \end{tabular}
  \end{threeparttable}
\end{table}

\subsection*{Conceptual Question: Search Rates vs.\ Hit Rates}

If groups with higher search rates also exhibit lower hit rates, this pattern
is consistent with lower suspicion thresholds being applied to those groups.
However, the outcome test is subject to the infra-marginality problem and
cannot detect discrimination occurring at the margin of search decisions.

\bibliographystyle{plainnat}
\bibliography{references}

\end{document}